\documentclass[10pt,a4paper]{amsart}
\usepackage[margin=19mm]{geometry}
\usepackage{amsmath,amssymb,mathtools}
\usepackage[scr=rsfs,cal=euler]{mathalfa}
\usepackage{xfrac}
\usepackage[unicode,hidelinks,bookmarks=false]{hyperref}
\newcommand{\ie}{i.e.\ }
\newcommand{\cf}{cf.\ }
\newcommand{\resp}{respectively\ }
\newcommand{\Subsection}{\subsection}
\providecommand{\nobreakdash}{\nobreak\hbox{-}}
\setlength{\emergencystretch}{2em}
\multlinegap0pt
\newtheorem{lemma}{Lemma}
\title{Bisections of graphs under degree constraints}
\author{Jie Ma, Hehui Wu}
\date{Source: arXiv:2504.15096v2 (CC BY 4.0)}
\begin{document}
\maketitle

\noindent\emph{Source and licence.} Bisections of graphs under degree constraints. Version arXiv:2504.15096v2 (CC BY 4.0), distributed by arXiv under the Creative Commons Attribution 4.0 International licence, http://creativecommons.org/licenses/by/4.0/, as recorded in the arXiv licence metadata for this exact version and shown on the abstract page https://arxiv.org/abs/2504.15096v2. Full licence notice: \texttt{licenses/arxiv-2504.15096-CC-BY-4.0.txt}.\par\smallskip

\noindent\textit{Editorial scope.} Counted unit: the article's lemma Lem:indegree-tripar2 (source0.tex lines 645--663). The article's complete original proof of it is delivered with it (source0.tex lines 665--815, one \texttt{proof} environment of 151 lines). No mathematics is written, repaired or re-derived: the statement and the proof are the article's own lines, in the article's own notation, and no retained line is substituted or re-flowed.  Retained uncounted as the source-local context the proof needs: lines 360--375; lines 533--535.  Nothing was omitted from any retained span, so the package carries no omission record.  No retained line contains a citation, so the wrapper carries no bibliography and no bibliography entry.  No retained line contains a figure environment, an image-inclusion command or drawing code, so no image is delivered and the package declares no asset dependency.  Editorial formatting only, in addition: an AMS \texttt{amsart} preamble that restates the article's own declarations and macros used by the retained lines, namely the environments \texttt{lemma} on separate counters, together with the standard packages those lines need.  The retained environments print in this document as Lemma 1 in order of appearance, whereas the article prints the same items with the numbering of its own full document.  The complete native source of the article as distributed in the arXiv e-print for this exact version is bundled unchanged as \texttt{sources/176848/source0.tex}.
	\begin{lemma}\label{Lem:key}
		Let $H$ be a graph with a family $\{A_i\}_{i \in I}$ of disjoint vertex subsets,\footnote{Here, $\bigcup_{i \in I} A_i$ may not equal $V(H)$.} where $I$ is a finite index set.
		For each $i \in I$, let $\eta_i > 0$ be a real number and let $a_i \geq 1$ be an integer. Define the subset
		\[
		A_i^+ = \{v \in A_i : d_H(v) \geq 2(1+\eta_i)a_i\}.
		\]
		and set $\eta = \min_{i \in I} \eta_i$. If
		\begin{equation}\label{equ:key}
			\left(1 + \frac{1}{\eta}\right) \sum_{i \in I} a_i |A_i \setminus A_i^+| < |V(H)|,
		\end{equation}
		then there exists a non-empty induced subgraph $H' \subseteq H$ such that
		\begin{itemize}
			\item [(a)] $d_{H'}(v) \geq a_i$ for all $v \in V(H') \cap A_i$ and $i\in I$;
			\item [(b)] $|V(H)\backslash V(H')| \leq \sum_{i \in I} a_i |A_i \setminus V(H')| \leq \left(1 + \frac{1}{\eta}\right) \sum_{i \in I} a_i |A_i \setminus A_i^+|$, where $V(H)\backslash V(H')\subseteq \bigcup_{i \in I} A_i$.
		\end{itemize}
	\end{lemma}

	\begin{equation}\label{equ:phi2}
		\varphi(i):=\varphi_{c,\varepsilon}(i)=\left(\frac{1-c}{4}-\frac{\mu_i}{2}\right)\cdot i, \mbox{~~where~~} \mu_i=4d_\varepsilon\cdot i^{\frac{1}{2}(\varepsilon-1)}+2\varepsilon.
	\end{equation}

	\begin{lemma}
		\label{Lem:indegree-tripar2}
		Let $A\cup B\cup C$ be a tripartition of $V(G)$ such that every vertex in $C$ is both $A$-good and $B$-good, $|A|>\frac{\varepsilon n}{1000}$, and
		\begin{equation}
			\label{equ:indegree condition2}
			\sum_{i\ge 1} i\cdot |(A\cap V^i)\backslash V^i_A|\le \frac{\varepsilon^2n}{250}.
		\end{equation}
		Then there exists a tripartition $A'\cup B'\cup C'$ of $V(G)$ such that the following hold:
		\begin{itemize}
			\item[(2a).] $B\subseteq B'$ and $C'\subseteq C$,\footnote{One cannot establish a containment relationship between $A$ and $A'$ in this lemma.}
			\item[(2b).] $|A|-\frac{\varepsilon n}{500}\le |A'|\le |A|+\frac{\varepsilon n}{500}$, $|B|\leq |B'|\leq |B|+\frac{\varepsilon n}{500}$ and $|C|-\frac{\varepsilon n}{500}\leq |C'|\leq |C|$,
			\item[(2c).] every vertex in $C'$ is both $A'$-good and $B'$-good,
			\item[(2d).] every vertex $x\in A'\cap V^i$ satisfies $d_{A'}(x)\ge \varphi(i)$, and
			\item[(2e).] every vertex in $B'\backslash B$ is $B'$-good. In particular, this shows that for each integer $i\geq 1$,
			\begin{equation}\label{equ:indeg B'}
				(B'\cap V^i)\backslash V^i_{B'}\subseteq (B\cap V^i)\backslash V^i_B.
			\end{equation}
		\end{itemize}
	\end{lemma}

	\begin{proof}
		Let $A\cup B\cup C$ be a tripartition of $V(G)$ given by the lemma.
		In what follows, beginning with $A\cup B\cup C$, we will define a sequence of operations that carefully relocate vertices.
		These operations will produce a series of tripartitions of $G$,
		each satisfying progressively stronger properties that bring us closer to the desired properties.
		Importantly, in each round, only a small fraction of vertices are affected.
		This ensures that the final tripartition retains a distribution very close to the original one.
		
		\medskip
		
		\noindent{\it Step 1: extracting a robust core inside $A$.}
		First, we identify a large subset $A^* \subseteq A$ satisfying property (2d).
		Later, we will see that this subset $A^*$ ensures that most vertices in $A$ remain unchanged throughout subsequent operations.
		To construct $A^*$, we apply Lemma~\ref{Lem:key} to $H=G[A]$.
		Let $I$ denote the set of all integers $1\leq i\leq n$ with $\varphi(i)>0$.
		Let $A_i:=A\cap V^i, \eta_i:=\mu_i$ and $a_i:=\varphi(i)$ for each $i\in I$.
		By the definition of $A_i^+$,
		we observe $$A_i^+ =\{v \in A_i : d_H(v) \geq 2(1+\eta_i)a_i\}=\{v\in A\cap V^i: d_A(v)\geq 2(1+\mu_i)\varphi(i)\}= A\cap V_A^i.$$
		Since $\eta=\min \eta_i\geq 2\varepsilon$ and $a_i=\varphi(i)\leq \frac{i}{4}$,
		it holds from \eqref{equ:indegree condition2} that
		$$\Omega:=(1+\frac{1}{\eta}) \sum_{i\in I} a_i \cdot|A_i\backslash A^+_i|\leq
		\frac{1}{\varepsilon} \sum_{i\in I} \frac{i}{4}\cdot |(A\cap V^i)\backslash V^i_A|\leq \frac{\varepsilon n}{1000}<|A|=|V(H)|,$$
		where the final inequality uses the assumption $|A|>\varepsilon n/1000$.
		By Lemma \ref{Lem:key}, there is a subset $A^*$ of $A$ such that
		\begin{equation}\label{equ:indeg deg(A*)}
			d_{A^*}(x)\ge \varphi(i) \text{~~for every } x\in A^*\cap V^i,
			\mbox{~~~and~~~}
			|A\backslash A^*|\le \sum_{i\in I} \varphi(i)\cdot |(A\cap V^i)\backslash A^*|\le \Omega\le \frac{\varepsilon n}{1000}.
		\end{equation}
		
		
		\medskip
		
		\noindent{\it Step 2: greedily removing bad vertices from $A$.}
		Next, we implement a greedy algorithm that iteratively moves a small proportion of vertices from $A \cup C$ into $B$.
		The algorithm produces an intermediate tripartition $V(G) = A_1 \cup B_1 \cup C_1$ satisfying $A^* \subseteq A_1$.
		Initially, let $(A^{0}, B^{0}, C^{0}) = (A, B, C)$.
		Given that $A^{k}\cup B^{k}\cup C^{k}$ is defined,
		if there exists a vertex $v_k\in A^{k}\cap V^i$ for some $i\in I$ satisfying
		\begin{equation}
			\label{equ:indeg alg}
			|N_G(v_k)\cap (A^{k}\cup C^{k})|<\varphi(i),
		\end{equation}
		then update $A^{k+1}\cup B^{k+1}\cup C^{k+1}$ by
		$$A^{k+1}=A^k-\{v_k\},~C^{k+1}=C^k\backslash N_G(v_k) \text{~~and~~} B^{k+1}=B^k\cup \{v_k\}\cup (N_G(v_k)\cap C^k);$$
		otherwise, terminate this algorithm and return the current tripartition as $A_1 \cup B_1 \cup C_1$.
		
		We analyze the properties of the tripartition $A_1 \cup B_1 \cup C_1$.
		By construction, for every iteration $k$ we have the following nested containments: $A_1\subseteq A^{k+1}\subseteq A^k\subseteq A$,
		$C_1\subseteq C^{k+1}\subseteq C^k\subseteq C$, and $B\subseteq B^k\subseteq B^{k+1}\subseteq B_1$.
		Moreover, every vertex $x\in A_1\cap V^i$ satisfies
		\begin{equation}
			\label{equ:indeg A'}
			|N_G(x)\cap (A_1\cup C_1)|\ge \varphi(i).
		\end{equation}
		Indeed, no vertex of $A^*$ can ever be selected as some $v_k$: if $x\in A^*\cap V^i$ and $x\in A^k$, then $A^*\subseteq A^k\cup C^k$ and hence
		$|N_G(x)\cap (A^k\cup C^k)|\ge d_{A^*}(x)\ge \varphi(i)$ by \eqref{equ:indeg deg(A*)}, contradicting \eqref{equ:indeg alg}.
		Therefore $A^*\subseteq A_1\subseteq A$, and
		\begin{equation}
			\label{equ:indeg A-A'}
			|A\backslash A_1|\le |A\backslash A^*|\le \frac{\varepsilon n}{1000},
		\end{equation}
		where $A\backslash A_1=\{v_k:k\ge 0\}$.
		Since each $v_k \in V^i$ satisfies $|N_G(v_k) \cap C^k| < \varphi(i)$ (by \eqref{equ:indeg alg}), we bound:
		\begin{equation}
			\label{equ:indeg C-C'}
			|C\backslash C_1|=\sum_{k} |N_G(v_k)\cap C^k|<\sum_{i\in I} \varphi(i)\cdot |(A\backslash A_1)\cap V^i|\le
			\sum_{i\in I} \varphi(i)\cdot |(A\backslash A^*)\cap V^i|\le \frac{\varepsilon n}{1000},
		\end{equation}
		where the last inequality follows from \eqref{equ:indeg deg(A*)}. Therefore,
		\begin{equation}
			\label{equ:indeg B'-B}
			|B_1\backslash B|=|A\backslash A_1|+|C\backslash C_1| \le \frac{\varepsilon n}{500}.
		\end{equation}
		
		We claim that for every $k\geq 0$,
		\begin{equation}
			\label{equ:indeg C^k}
			\text{ each vertex in } C^k \text{ is both } A^k\text{-good  and } B^k\text{-good}.
		\end{equation}
		We prove this by induction on $k$. The base case $k=0$ holds since every vertex in $C$ is $A$-good and $B$-good by assumption.
		Now suppose \eqref{equ:indeg C^k} holds for some $k \geq 0$. Consider any vertex $x \in C^{k+1} = C^k \setminus N_G(v_k)$.
		Since $x\in C^{k+1}\subseteq C^k$, by induction $x$ is $B^k$-good.
		Since $B^k \subseteq B^{k+1}$, this evidently implies that $x$ is also $B^{k+1}$-good.
		Moreover, since $x$ is not adjacent to $v_k$, crucially we have $d_{A^k}(x) = d_{A^{k+1}}(x)$,
		which implies that $x$ remains $A^{k+1}$-good.
		This proves \eqref{equ:indeg C^k} and in particular, implies that
		\begin{equation}
			\label{equ:indeg C'}
			\text{ each vertex in } C_1 \text{ is both } A_1\text{-good  and } B_1\text{-good}.
		\end{equation}
		
		We also claim that for every $k\geq 0$,
		\begin{equation}
			\label{equ:indeg B^k}
			\text{ each vertex in } B^{k+1}\backslash B^k \text{ is } B^k\text{-good}.
		\end{equation}
		Note that $B^{k+1}\backslash B^k =\{v_k\}\cup (N_G(v_k)\cap C^k)$, so by \eqref{equ:indeg C^k}
		it is enough to show that $v_k$ is $B^k$-good, i.e., $|N_{G}(v_k)\cap B^k|\geq 2(1+\mu_i)\varphi(i)$ when $v_k\in V^i$.
		By \eqref{equ:indeg alg}, we have $|N_G(v_k)\cap B^k|=i-|N_G(v_k)\cap (A^k\cup C^k)|> i-\varphi(i)$.
		Using \eqref{equ:phi2}, we compute
		\begin{align*}
			i-\varphi(i)-2(1+\mu_i)\varphi(i)
			&=\left(1-(3+2\mu_i)\left(\frac{1-c}{4}-\frac{\mu_i}{2}\right)\right)i\\
			&=\left(\frac{1+3c}{4}+\left(1+\frac{c}{2}\right)\mu_i+\mu_i^2\right)i\ge 0.
		\end{align*}
		Hence $i-\varphi(i)\ge 2(1+\mu_i)\varphi(i)$, proving \eqref{equ:indeg B^k}.
		Since $B^k$ grows to $B_1$ as $k$ increases, we derive from \eqref{equ:indeg B^k} that
		\begin{equation}
			\label{equ:indeg B-good}
			\text{ each vertex in } B_1\backslash B \text{ is } B_1\text{-good}.
		\end{equation}
		
		
		\medskip
		
		\noindent{\it Step 3: repairing the remaining deficits in $A_1$.}
		Let us point out that the tripartition $A_1 \cup B_1 \cup C_1$ satisfies all desired properties of Lemma~\ref{Lem:indegree-tripar2}, except for property (2d).
		To fix this, we move some vertices from $C_1$ into $A_1$ as follows.
		For every $x\in A_1\cap V^i$,
		if $d_{A_1}(x)<\varphi(i)$, then in view of \eqref{equ:indeg A'} there exists a subset $R_x$
		containing exactly $\varphi(i)-d_{A_1}(x)$ neighbors of $x$ in $C_1$;
		otherwise $d_{A_1}(x)\geq \varphi(i)$, let $R_x=\emptyset$.
		Note that if $R_x\neq \emptyset$ for some $x\in A_1$, then $d_{A^*}(x)\le d_{A_1}(x)<\varphi(i)$,
		which together with \eqref{equ:indeg deg(A*)} imply that $x\in A_1\backslash A^*$.
		We then move all vertices in the set
		\begin{equation}\label{equ:indeg R}
			R:=\cup_{x\in A_1} R_x=\cup_{x\in A_1\backslash A^*} R_x
		\end{equation}
		from $C_1$ into $A_1$.
		This results in a new tripartition, denoted by $V(G)=A_2\cup B_2\cup C_2$, where
		$$A_2=A_1\cup R,~ B_2=B_1, \text{ and } C_2=C_1\backslash R.$$
		
		We now show $A_2\cup B_2\cup C_2$ is the desired tripartition for this lemma.
		Clearly, $B\subseteq B_1=B_2$ and $C_2\subseteq C_1\subseteq C$, thus property (2a) holds.
		For every $x\in A_1\cap V^i$, from the construction we see $d_{A_2}(x)\geq \varphi(i)$.
		By \eqref{equ:indeg C'}, every vertex in $R\subseteq C_1$ is $A_1$-good.
		Therefore, every vertex $x\in A_2\cap V^i$ satisfies that $d_{A_2}(x)\ge \varphi(i)$,
		establishing property (2d) for $A_2\cup B_2\cup C_2$.
		Since $A_1\subseteq A_2$ and $B_1=B_2$,
		every $A_1$-good (or $B_1$-good) vertex remains $A_2$-good (or $B_2$-good).
		Therefore \eqref{equ:indeg C'} and \eqref{equ:indeg B-good} imply properties (2c) and (2e) for $A_2\cup B_2\cup C_2$, respectively.
		It remains to show property (2b) for $A_2\cup B_2\cup C_2$.
		To do so, we need to estimate the size of $R$.
		We can derive that
		$$|R|\le \sum_{x\in A_1\backslash A^*} |R_x|\le \sum_{i} \varphi(i)\cdot |(A\cap V^i)\backslash A^*|\le \frac{\varepsilon n}{1000},$$
		where the first inequality follows by \eqref{equ:indeg R}, the second one follows by the fact $A_1\subseteq A$ and by the construction that $|R_x|\leq \varphi(i)$ for each $x\in A_1\cap V^i$,
		and the last one follows by \eqref{equ:indeg deg(A*)}.
		This upper bound on $|R|$, together with \eqref{equ:indeg A-A'}, \eqref{equ:indeg C-C'}, and \eqref{equ:indeg B'-B},
		establishes property (2b) for the tripartition $A_2 \cup B_2 \cup C_2$, thereby completing the proof of Lemma~\ref{Lem:indegree-tripar2}.
	\end{proof}

\end{document}
